\setcounter{section}{3}
\setlength{\parindent}{0pt}  % niente rientro
\setlength{\parskip}{6pt}
\setlength{\itemsep}{3pt}
\section{User Story}
\subsection{User Story legate ai requisiti funzionali comuni a tutti gli agenti}
\subsubsection*{User Story 1 – associata a RF-01: Registrazione utente}
\textbf{Registrazione autonoma con validazione dei dati}

Come utente, voglio potermi registrare autonomamente alla webapp inserendo i miei dati personali, in modo da poter creare un profilo e accedere ai servizi della piattaforma.

\textbf{Criteri di accettazione:}
\begin{itemize}
    \item Tutti i campi obbligatori devono essere compilati (nome, cognome, email, codice fiscale, numero di telefono, comune di residenza, password).
    \item Il sistema valida i formati (email, CF, telefono, password) e verifica l'unicità dell'email e del CF.
    \item In caso di errori, vengono mostrati messaggi specifici.
    \item Dopo la conferma, il sistema crea l’account e invia un’email di verifica.
\end{itemize}
\textbf{Tasks:}
\begin{itemize}
    \item Creare il form di registrazione con validazione client-side e server-side.
    \item Implementare i controlli di formato per ciascun campo.
    \item Integrare l’invio email di conferma registrazione.
    \item Gestire messaggi di errore e conferma in UI.
\end{itemize}
\subsubsection*{User Story 2 – associata a RF-02: Login utente}
\textbf{Accesso con credenziali e gestione errori}

Come utente registrato, voglio poter accedere alla webapp inserendo le mie credenziali, in modo da poter visualizzare la mia area personale.

\textbf{Criteri di accettazione:}
\begin{itemize}
    \item L'utente inserisce email e password valide.
    \item Il sistema autentica e reindirizza alla dashboard personale.
    \item In caso di errore, viene mostrato un messaggio (“Credenziali errate”).
\end{itemize}
\textbf{Tasks:}
\begin{itemize}
    \item Implementare il form di login con controlli di validazione.
    \item Creare la logica di autenticazione e gestione sessione.
    \item Implementare la visualizzazione dei messaggi di errore nel form di login.
    \item Gestire il reindirizzamento post-login.
\end{itemize}
\subsubsection*{User Story 3 – associata a RF-02.1: Recupero password}
\textbf{Reimpostazione password tramite email}

Come utente, voglio poter recuperare la mia password se la dimentico, in modo da poter accedere di nuovo al mio account.

\textbf{Criteri di accettazione:}
\begin{itemize}
    \item L’utente inserisce la propria email.
    \item Il sistema invia un link sicuro per il reset.
    \item Il link reindirizza a una pagina di nuova password conforme ai criteri di sicurezza.
\end{itemize}
\textbf{Tasks:}
\begin{itemize}
    \item Creare la pagina “Password dimenticata”.
    \item Implementare la generazione di token di reset e scadenza temporale.
    \item Inviare email con link personalizzato.
    \item Validare e aggiornare la nuova password.
\end{itemize}
\subsubsection*{User Story 4 – associata a RF-03: Accesso all’area personale}
\textbf{Gestione area personale utente autenticato}

Come utente autenticato, voglio accedere alla mia area personale per visualizzare e gestire i miei dati.

\textbf{Criteri di accettazione:}
\begin{itemize}
    \item L’area è accessibile solo se l’utente è autenticato.
    \item Sono visibili dati anagrafici e impostazioni account.
    \item Accesso protetto tramite sessione attiva.
\end{itemize}
\textbf{Tasks:}
\begin{itemize}
    \item Implementare routing protetto per l’area personale.
    \item Collegare la visualizzazione dei dati utente dal database.
    \item Gestire errori di accesso non autorizzato.
\end{itemize}
\subsubsection*{User Story 5 – associata a RF-03.1-03.4: Gestione dati profilo e account}
\textbf{Aggiornamento o cancellazione delle informazioni personali}

Come utente autenticato, voglio poter modificare o cancellare le mie informazioni personali, in modo da mantenere aggiornati i miei dati e gestire la mia privacy.

\textbf{Criteri di accettazione:}
\begin{itemize}
    \item L’utente può aggiornare email, telefono, dati non obbligatori.
    \item Ogni modifica è validata e confermata.
    \item In caso di cancellazione account, viene chiesta conferma e rimossi i dati.
\end{itemize}
\textbf{Tasks:}
\begin{itemize}
    \item Creare form di modifica profilo.
    \item Validare i nuovi dati (formato email, telefono).
    \item Implementare procedura di cancellazione con conferma.
    \item Inviare notifica email dopo aggiornamento o eliminazione.
\end{itemize}
\subsubsection*{User Story 6 – associata a RF-04: Logout utente}
\textbf{Logout sicuro dalla webapp}

Come utente autenticato, voglio potermi disconnettere in modo sicuro, in modo da impedire accessi non autorizzati.

\textbf{Criteri di accettazione:}
\begin{itemize}
    \item Alla disconnessione, la sessione viene invalidata.
    \item L’utente è reindirizzato alla homepage pubblica.
\end{itemize}
\textbf{Tasks:}
\begin{itemize}
    \item Implementare funzione di logout nel menu utente.
    \item Cancellare token/sessione al logout.
    \item Gestire redirect alla pagina iniziale.
\end{itemize}
\subsection{User Story legate ai requisiti funzionali specifici per gli amministratori della CER}
\subsubsection*{User Story 7 – associata a RF-05: Gestione membri della CER}
\textbf{Gestione elenco membri e loro dati}

Come amministratore della CER, voglio visualizzare e gestire i membri della comunità, in modo da poter mantenere aggiornata la composizione della CER.

\textbf{Criteri di accettazione:}
\begin{itemize}
    \item L’amministratore può aggiungere, modificare o rimuovere membri.
    \item Ogni operazione è tracciata.
    \item Lo stato di ciascun membro è visibile (attivo, sospeso, in attesa).
\end{itemize}
\textbf{Tasks:}
\begin{itemize}
    \item Creare tabella membri con azioni CRUD.
    \item Implementare log delle modifiche.
    \item Aggiornare stato membro in tempo reale.
\end{itemize}
\subsubsection*{User Story 8 – associata a RF-06: Approvazione richieste di adesione}
\textbf{Gestione richieste di nuovi membri}

Come amministratore della CER, voglio approvare o rifiutare le richieste di adesione, in modo da controllare chi entra nella comunità.

\textbf{Criteri di accettazione:}
\begin{itemize}
    \item Le richieste sono visibili in elenco.
    \item L’amministratore può accettare o rifiutare con motivazione.
    \item L’utente riceve una notifica email del risultato.
\end{itemize}
\textbf{Tasks:}
\begin{itemize}
    \item Implementare sezione “Richieste in attesa”.
    \item Creare azioni di approvazione/rifiuto.
    \item Collegare invio automatico di notifica email.
\end{itemize}
\subsubsection*{User Story 9 – associata a RF-07-08: Gestione e monitoraggio impianti}
\textbf{Gestione impianti e monitoraggio stato operativo}

Come amministratore della CER, voglio registrare e monitorare gli impianti energetici, in modo da avere una visione completa del funzionamento della comunità.

\textbf{Criteri di accettazione:}
\begin{itemize}
    \item Possibile aggiungere/modificare impianti con dati tecnici.
    \item Visualizzazione in tabella e/o mappa.
    \item Stato aggiornato (attivo, guasto, manutenzione).
\end{itemize}
\textbf{Tasks:}
\begin{itemize}
    \item Creare modulo di registrazione impianto.
    \item Implementare monitoraggio con aggiornamento automatico.
    \item Mostrare alert in caso di guasto.
\end{itemize}
\subsubsection*{User Story 10 – associata a RF-09-10: Reportistica e analisi energetica}
\textbf{Visualizzazione report e analisi dei benefici}

Come amministratore della CER, voglio consultare report energetici e analisi economiche, in modo da monitorare l’efficienza e i benefici della CER.

\textbf{Criteri di accettazione:}
\begin{itemize}
    \item Possibilità di selezionare periodo temporale.
    \item Dati mostrati in grafici e tabelle.
    \item Esportazione disponibile in CSV o PDF.
\end{itemize}
\textbf{Tasks:}
\begin{itemize}
    \item Implementare dashboard analitica.
    \item Collegare grafici interattivi a database.
    \item Creare funzioni di esportazione.
\end{itemize}
\subsection{User story legate ai requisti funzionali specifici per il Rappresentante del Comune}
\subsubsection*{User Story 11 – associata a RF-11-12: Supervisione e dashboard comunale}
\textbf{Visualizzazione delle CER del territorio}

Come rappresentante del Comune, voglio visualizzare le CER attive e i dati aggregati, in modo da monitorare la produzione e l’impatto energetico sul territorio.

\textbf{Criteri di accettazione:}
\begin{itemize}
    \item Possibile vedere elenco CER con indicatori principali.
    \item Dashboard aggregata per periodo e territorio.
    \item Nessuna modifica diretta sui dati delle CER.
\end{itemize}
\textbf{Tasks:}
\begin{itemize}
    \item Creare vista “Panoramica comunale”.
    \item Implementare filtri temporali e geografici.
    \item Collegare grafici ai dati aggregati.
\end{itemize}
\subsubsection*{User Story 12 – associata a RF-13: Gestione richieste dalle CER}
\textbf{Supervisione delle richieste di supporto}

Come rappresentante del Comune, voglio gestire le richieste provenienti dalle CER, in modo da fornire supporto o approvazioni necessarie.

\textbf{Criteri di accettazione:}
\begin{itemize}
    \item Visualizzazione delle richieste in ordine cronologico.
    \item Possibilità di modificare stato (in lavorazione, approvata, rifiutata).
    \item Possibilità di aggiungere un commento di risposta.
\end{itemize}
\textbf{Tasks:}
\begin{itemize}
    \item Creare sezione “Richieste CER”.
    \item Implementare pulsanti per aggiornare stato.
    \item Gestire notifiche agli amministratori CER.
\end{itemize}
\subsubsection*{User Story 13 – associata a RF-14: Comunicazioni istituzionali}
\textbf{Invio comunicazioni alle CER del Comune}

Come rappresentante del Comune, voglio poter inviare comunicazioni ufficiali alle CER registrate, in modo da informarle su iniziative e novità.

\textbf{Criteri di accettazione:}
\begin{itemize}
    \item Possibilità di scrivere testo e allegare file.
    \item Le comunicazioni vengono recapitate nell’area notifiche CER.
\end{itemize}
\textbf{Tasks:}
\begin{itemize}
    \item Implementare modulo di invio messaggi.
    \item Gestire upload e invio allegati.
    \item Creare visualizzazione comunicazioni in area notifiche.
\end{itemize}
\subsection{User story legate ai requisti funzionali specifici per i membri della CER}
\subsubsection*{User Story 14 – associata a RF-15: Visualizzazione dati energetici personali}
\textbf{Monitoraggio dei propri consumi e produzione}

Come membro della CER, voglio visualizzare i miei dati energetici personali, in modo da capire i miei consumi e la quota condivisa di energia.

\textbf{Criteri di accettazione:}
\begin{itemize}
    \item Dashboard con grafici dei dati energetici.
    \item Possibilità di filtrare per periodo.
    \item Aggiornamento automatico o periodico.
\end{itemize}
\textbf{Tasks:}
\begin{itemize}
    \item Creare dashboard personale.
    \item Collegare API dati energetici.
    \item Implementare filtri temporali e refresh automatico.
\end{itemize}
\subsubsection*{User Story 15 – associata a RF-16: Confronto con la comunità}
\textbf{Confronto dei propri dati con la media della CER}

Come membro della CER, voglio confrontare i miei consumi con la media della CER, in modo da valutare la mia efficienza energetica.

\textbf{Criteri di accettazione:}
\begin{itemize}
    \item Mostra grafico comparativo utente vs media CER.
    \item Dati anonimi e aggregati.
\end{itemize}
\textbf{Tasks:}
\begin{itemize}
    \item Implementare logica di calcolo media comunitaria.
    \item Visualizzare confronto grafico in dashboard.
\end{itemize}
\subsubsection*{User Story 16 – associata a RF-17: Notifiche e avvisi}
\textbf{Ricezione di notifiche e aggiornamenti}

Come membro della CER, voglio ricevere notifiche sulle attività della mia CER, in modo da essere sempre informato su cambiamenti e novità.

\textbf{Criteri di accettazione:}
\begin{itemize}
    \item Le notifiche appaiono in area dedicata e sono marcabili come lette.
    \item Devono coprire eventi come approvazioni, anomalie o nuovi impianti.
\end{itemize}
\textbf{Tasks:}
\begin{itemize}
    \item Implementare centro notifiche.
    \item Gestire logica di creazione e lettura notifiche.
    \item Creare contatore visivo di notifiche non lette.
\end{itemize}
\subsubsection*{User Story 17 – associata a RF-18: Segnalazioni e richieste}
\textbf{Invio segnalazioni tecniche o amministrative}

Come membro della CER, voglio poter inviare segnalazioni alla CER o al Comune, in modo da comunicare problemi o suggerimenti.

\textbf{Criteri di accettazione:}
\begin{itemize}
    \item Modulo con titolo, descrizione e allegati.
    \item Tracciamento stato segnalazione (inviata, in lavorazione, risolta).
\end{itemize}
\textbf{Tasks:}
\begin{itemize}
    \item Creare modulo segnalazione con upload file.
    \item Gestire backend per stato della richiesta.
    \item Inviare notifiche automatiche ai destinatari.
\end{itemize}